\documentclass[10pt,a4paper]{amsart}
\usepackage{amsmath,amssymb,amsthm,graphicx}
% Retained preamble for the Stacks Project chapter snapshot in this directory.
% The chapter's own preamble.tex is a copy of the shared Stacks Project preamble; the block
% below is sliced MECHANICALLY, line for line, from it. Nothing was added, removed or
% redefined; the wrapper class, the xr-hyper/hyperref wiring and the externaldocument list
% are omitted because the assembler supplies the document class.
%
% lines taken from preamble.tex: 3, 5, 9, 11, 12, 156, 157, 158, ... (51 lines total)
%
% package coverage restored from the chapter's own preamble (2026-10-13): xy, xy-2cell

\usepackage[all]{xy}
\xyoption{2cell}
\UseAllTwocells

\newenvironment{reference}{\comment}{\endcomment}
\newenvironment{slogan}{\comment}{\endcomment}
\newenvironment{history}{\comment}{\endcomment}
\theoremstyle{plain}
\newtheorem{theorem}[subsection]{Theorem}
\newtheorem{proposition}[subsection]{Proposition}
\newtheorem{lemma}[subsection]{Lemma}
\theoremstyle{definition}
\newtheorem{definition}[subsection]{Definition}
\newtheorem{example}[subsection]{Example}
\newtheorem{exercise}[subsection]{Exercise}
\newtheorem{situation}[subsection]{Situation}
\theoremstyle{remark}
\newtheorem{remark}[subsection]{Remark}
\newtheorem{remarks}[subsection]{Remarks}
\numberwithin{equation}{subsection}
\def\lim{\mathop{\mathrm{lim}}\nolimits}
\def\colim{\mathop{\mathrm{colim}}\nolimits}
\def\Spec{\mathop{\mathrm{Spec}}}
\def\Hom{\mathop{\mathrm{Hom}}\nolimits}
\def\Ext{\mathop{\mathrm{Ext}}\nolimits}
\def\SheafHom{\mathop{\mathcal{H}\!\mathit{om}}\nolimits}
\def\SheafExt{\mathop{\mathcal{E}\!\mathit{xt}}\nolimits}
\def\Sch{\mathit{Sch}}
\def\Mor{\mathop{\mathrm{Mor}}\nolimits}
\def\Ob{\mathop{\mathrm{Ob}}\nolimits}
\def\Sh{\mathop{\mathit{Sh}}\nolimits}
\def\NL{\mathop{N\!L}\nolimits}
\def\CH{\mathop{\mathrm{CH}}\nolimits}
\def\proetale{{pro\text{-}\acute{e}tale}}
\def\etale{{\acute{e}tale}}
\def\QCoh{\mathit{QCoh}}
\def\Ker{\mathop{\mathrm{Ker}}}
\def\Im{\mathop{\mathrm{Im}}}
\def\Coker{\mathop{\mathrm{Coker}}}
\def\Coim{\mathop{\mathrm{Coim}}}
\DeclareMathSymbol{\boxtimes}{\mathbin}{AMSa}{"02}
\def\QCohstack{\mathcal{QC}\!\mathit{oh}}
\def\Cohstack{\mathcal{C}\!\mathit{oh}}
\def\Spacesstack{\mathcal{S}\!\mathit{paces}}
\def\Quotfunctor{\mathrm{Quot}}
\def\Hilbfunctor{\mathrm{Hilb}}
\def\Curvesstack{\mathcal{C}\!\mathit{urves}}
\def\Polarizedstack{\mathcal{P}\!\mathit{olarized}}
\def\Complexesstack{\mathcal{C}\!\mathit{omplexes}}
\def\Pic{\mathop{\mathrm{Pic}}\nolimits}
\def\Picardstack{\mathcal{P}\!\mathit{ic}}
\def\Picardfunctor{\mathrm{Pic}}
\def\Deformationcategory{\mathcal{D}\!\mathit{ef}}
\usepackage{cleveref}
\crefname{originallemma}{original lemma}{original lemmas}
\Crefname{originallemma}{Original lemma}{Original lemmas}
\makeatletter
\newlabel{etale-cohomology-lemma-compare-cohomology}{{0DDH}{source}{}{sourcepointer-1}{}}
\newlabel{etale-cohomology-lemma-compare-cohomology@cref}{{[originallemma][0][]0DDH}{[1][0][]source}{}{}{}}
\makeatother
\crefname{originallemma}{original lemma}{original lemmas}
\Crefname{originallemma}{Original lemma}{Original lemmas}
\makeatletter
\newlabel{perfect-lemma-affine-compare-bounded}{{06Z0}{source}{}{sourcepointer-2}{}}
\newlabel{perfect-lemma-affine-compare-bounded@cref}{{[originallemma][0][]06Z0}{[1][0][]source}{}{}{}}
\makeatother
\crefname{originallemma}{original lemma}{original lemmas}
\Crefname{originallemma}{Original lemma}{Original lemmas}
\makeatletter
\newlabel{section-restriction}{{section-restriction}{source}{}{sourcepointer-3}{}}
\newlabel{section-restriction@cref}{{[originallemma][0][]section-restriction}{[1][0][]source}{}{}{}}
\makeatother
\crefname{originallemma}{original lemma}{original lemmas}
\Crefname{originallemma}{Original lemma}{Original lemmas}
\makeatletter
\newlabel{sites-cohomology-lemma-cartesion-plus-topology}{{0GZS}{source}{}{sourcepointer-4}{}}
\newlabel{sites-cohomology-lemma-cartesion-plus-topology@cref}{{[originallemma][0][]0GZS}{[1][0][]source}{}{}{}}
\makeatother
\crefname{originallemma}{original lemma}{original lemmas}
\Crefname{originallemma}{Original lemma}{Original lemmas}
\makeatletter
\newlabel{spaces-perfect-remark-match-total-direct-images}{{08GH}{source}{}{sourcepointer-5}{}}
\newlabel{spaces-perfect-remark-match-total-direct-images@cref}{{[originallemma][0][]08GH}{[1][0][]source}{}{}{}}
\makeatother
\title{Stacks, Sheaves on Algebraic Stacks: lemma-QC-compare}
\author{The Stacks Project Authors}
\date{stacks-project:stacks-sheaves.tex:local-snapshot-20261012; GNU FDL 1.2 or later}
\begin{document}
\maketitle
\noindent\textit{Editorial note.} Statement span is lines 4961-4968 of sources/stacks-algebraic-stacks-sheaves/pinned-original-source.tex; proof span is lines 4970-4994, the complete author proof (23 body lines).\par\medskip
\noindent\textit{Editorial note.} Navigation stubs: 5 cited-but-not-included labels are supplied as non-load-bearing pointers with their pinned tags where the snapshot tag index has them.\par\medskip
\noindent\textit{Editorial note.} extra\_packages is empty on purpose: the retained preamble already loads xy with options and a second declaration clashes in the assembler wrapper.\par\medskip
\noindent\textit{Editorial note.} Text is retained verbatim, including the author's own typos and placeholders.\par\medskip
\noindent\textit{Editorial note.} Queue index 21 of queue\_algebraic-stacks-sheaves.json (pb=1217, stmt 4961-4968, proof 4970-4994).\par\medskip
\noindent\textit{Allowed standard prerequisites (editorial).}\begin{itemize}
\item Explicit prerequisite boundary. The cited inputs are prerequisites; the argument in the retained span is not replaced by any cited result.
\end{itemize}\medskip
\ifdefined\hypertarget\hypertarget{sourcepointer-1}{}\fi
\noindent\textit{Source reference.} Original lemma 0DDH: sources/stacks-algebraic-stacks-sheaves/pinned-original-source.tex, source lines 4994--4994. Cited by the retained statement or proof but not part of the collected spans; supplied only as a navigation pointer. The argument does not depend on its content.\par\medskip
\ifdefined\hypertarget\hypertarget{sourcepointer-2}{}\fi
\noindent\textit{Source reference.} Original lemma 06Z0: sources/stacks-algebraic-stacks-sheaves/pinned-original-source.tex, source lines 4994--4994. Cited by the retained statement or proof but not part of the collected spans; supplied only as a navigation pointer. The argument does not depend on its content.\par\medskip
\ifdefined\hypertarget\hypertarget{sourcepointer-3}{}\fi
\noindent\textit{Source reference.} Original lemma section-restriction: sources/stacks-algebraic-stacks-sheaves/pinned-original-source.tex, source lines 1185--1200. Cited by the retained statement or proof but not part of the collected spans; supplied only as a navigation pointer. The argument does not depend on its content.\par\medskip
\ifdefined\hypertarget\hypertarget{sourcepointer-4}{}\fi
\noindent\textit{Source reference.} Original lemma 0GZS: sources/stacks-algebraic-stacks-sheaves/pinned-original-source.tex, source lines 4994--4994. Cited by the retained statement or proof but not part of the collected spans; supplied only as a navigation pointer. The argument does not depend on its content.\par\medskip
\ifdefined\hypertarget\hypertarget{sourcepointer-5}{}\fi
\noindent\textit{Source reference.} Original lemma 08GH: sources/stacks-algebraic-stacks-sheaves/pinned-original-source.tex, source lines 4994--4994. Cited by the retained statement or proof but not part of the collected spans; supplied only as a navigation pointer. The argument does not depend on its content.\par\medskip
\begin{lemma}
\label{lemma-QC-compare}
Let $S$ be a scheme. Let $\mathcal{X} \to (\Sch/S)_{fppf}$ be a category
fibred in groupoids. The comparison morphism
$\epsilon : \mathcal{X}_{affine, \etale} \to \mathcal{X}_{affine}$
satisfies the assumptions and conclusions of Cohomology on Sites, Lemma
\ref{sites-cohomology-lemma-cartesion-plus-topology}.
\end{lemma}

\begin{proof}
Assumption (1) holds by definition of $\mathcal{X}_{affine}$.
For condition (2) we use that for $x \in \Ob(\mathcal{X})$
lying over the affine scheme $U = p(x)$ we have an equivalence
$\mathcal{X}_{affine, \etale}/x = (\textit{Aff}/U)_\etale$ compatible
with structure sheaves; see
discussion in Section \ref{section-restriction}.
Thus it suffices to show: given an affine scheme $U = \Spec(R)$
and a complex of $R$-modules $M^\bullet$ the total cohomology
of the complex of modules on $(\textit{Aff}/U)_\etale$
associated to $M^\bullet$ is quasi-isomorphic to $M^\bullet$.
This follows from a combination of:
Derived Categories of Schemes, Lemma
\ref{perfect-lemma-affine-compare-bounded} (total cohomology of complexes
of modules over affines in the Zariski topology),
Derived Categories of Spaces, Remark
\ref{spaces-perfect-remark-match-total-direct-images}
(agreement between total cohomology in small
Zariski and \'etale topologies for
quasi-coherent complexes of modules), and
\'Etale Cohomology, Lemma \ref{etale-cohomology-lemma-compare-cohomology}
(to see that the \'etale cohomology of a complex of modules
on the big \'etale site of a scheme
may be computed after restricting to the small \'etale site).
\end{proof}

\end{document}
